% =============================================================================
% Tianle Zhang — Curriculum Vitae
% Compile with XeLaTeX. Auto-built by GitHub Actions (.github/workflows/build-cv.yml).
% =============================================================================
\documentclass[10pt,a4paper]{article}

% ---------- geometry ----------
\usepackage[a4paper,top=1.4cm,bottom=1.4cm,left=1.6cm,right=1.6cm]{geometry}

% ---------- fonts (sans-serif body, close to Calibri feel) ----------
\usepackage{fontspec}
\setmainfont{TeX Gyre Heros}            % Helvetica-like sans
\newfontfamily\cjkfont{Droid Sans Fallback}  % fallback for "张天乐" locally; GA will use Noto
\linespread{1.05}

% ---------- packages ----------
\usepackage{titlesec}
\usepackage{enumitem}
\usepackage{xcolor}
\usepackage[hidelinks]{hyperref}

% ---------- colours / links ----------
\definecolor{accent}{HTML}{2C5282}
\hypersetup{colorlinks=true, urlcolor=accent, linkcolor=accent}

% ---------- section style: small caps + thin rule, tight ----------
\titleformat{\section}
  {\normalfont\large\bfseries\scshape}{}{0em}{}[\vspace{-4pt}\noindent\textcolor{black!40}{\rule{\linewidth}{0.4pt}}]
\titlespacing*{\section}{0pt}{8pt}{2pt}

% ---------- helpers ----------
\newcommand{\me}[1]{\textbf{#1}}
\newcommand{\eq}{\textsuperscript{\textbf{*}}}
\newcommand{\corr}{\textsuperscript{\textbf{\dag}}}

% Two-column row: bold left + grey-ish right, on one line
\newcommand{\row}[2]{\noindent\textbf{#1}\hfill #2\par}
\newcommand{\subrow}[2]{\noindent\textit{#1}\hfill #2\par}

% ---------- list defaults: tight ----------
\setlist[itemize]{leftmargin=1.2em, topsep=1pt, itemsep=1pt, partopsep=0pt, parsep=0pt, label=\textbullet}
\setlist[enumerate]{leftmargin=1.5em, topsep=2pt, itemsep=3pt, partopsep=0pt, parsep=0pt}

\setlength{\parindent}{0pt}
\setlength{\parskip}{1pt}
\pagestyle{empty}

% =============================================================================
\begin{document}

% ---------- HEADER ----------
{\Huge\bfseries Tianle Zhang \,\,{\normalsize\mdseries\,({\cjkfont 张天乐})}}\par
\vspace{2pt}
Ph.D.\ Student, College of Computing and Data Science, Nanyang Technological University, Singapore\par
\vspace{1pt}
\href{mailto:tianle006@e.ntu.edu.sg}{tianle006@e.ntu.edu.sg} \,|\, +65 8491 7566 \,|\,
\href{https://ztlmememe.github.io/}{Homepage} \,|\,
\href{https://scholar.google.com.hk/citations?user=Qbd4ybsAAAAJ&hl=en}{Google Scholar} \,|\, Male\par
\vspace{2pt}
{\small\textbf{Research interests:}\;
\#Multi-modal LLMs \,$\cdot$\, \#AI4S \,$\cdot$\, \#Efficient Deep Learning \,$\cdot$\, \#Evaluation of Generative Models \,$\cdot$\, \#Dataset Distillation}

% ---------- PROFILE ----------
\section*{Profile}
{\small
My research focuses on combining large language models (LLMs) with domain foundation models (FMs) ---
making the integration both efficient and effective --- with growing emphasis on scientific domains
such as small molecules and proteins. On the modeling side, I proposed a training-free framework
that enables LLMs to reason over non-text modalities via in-context representation learning, with
the molecular domain as the testbed (NeurIPS 2025, first author; Prof.\ Alvin Chan as corresponding
author), and contributed to subsequent work on strategic multimodal representation alignment
(AAAI 2026). On the evaluation side, I re-designed the human-evaluation protocol for text-to-video
generation (NeurIPS 2024, first author), and contributed to ConvBench, a hierarchical multi-turn
benchmark for large vision-language models (NeurIPS 2024 Spotlight). Earlier work on efficient
learning includes expanding-window trajectory matching for lossless graph dataset condensation
(ICML 2024, co-first author). Across 8 publications at top-tier ML venues (NeurIPS, ICML, AAAI) and
IEEE conferences, I am first or co-first author on 5. I am supported by the NTU Research
Scholarship, and serve as a reviewer for ICLR 2026, NeurIPS 2025/2026, and ICML 2025/2026.
Further details on my \href{https://ztlmememe.github.io/}{homepage}.
}

% ---------- EDUCATION ----------
\section*{Education}

\row{Nanyang Technological University, College of Computing and Data Science}{Singapore \,|\, Jan.~2025 -- Present}
Ph.D.\ in Computer Science;\, GPA: 5.0/5.0. \,
Advisor: \textit{Prof.\ Alvin Chan Guo Wei}. \,
Funded by \textit{NTU Research Scholarship} (full tuition + stipend, Sep.~2024).

\vspace{4pt}
\row{University of Electronic Science and Technology of China}{Chengdu, China \,|\, Sep.~2020 -- Jun.~2024}
B.Eng.\ in Computer Science and Technology (Joint Degree in Smart Finance \& Blockchain Finance); GPA: 88.51/100.\par
{\small\textit{Selected coursework:} Machine Learning (97), Data Structures \& Algorithms (95), Smart Contract \& DApps (95), Mathematical Modeling (96), Data Mining \& Big Data (92), Probability \& Statistics (92), Database Theory (91), Software Engineering (91).}

% ---------- ACADEMIC EXPERIENCE ----------
\section*{Academic Experience}

\row{Nanyang Technological University, CCDS}{Singapore}
\subrow{Ph.D.\ Student.\, Advisor: Prof.\ Alvin Chan Guo Wei}{Jan.~2025 -- Present}
\begin{itemize}
  \item Exploring new ways to combine LLMs with domain foundation models (FMs), aiming for integrations that are both more efficient and more effective. Current focus: scientific domains --- small molecules and proteins --- where structured FM representations complement LLM reasoning.
\end{itemize}

\row{Nanyang Technological University, CCDS}{Singapore}
\subrow{Research Assistant.\, Advisor: Prof.\ Alvin Chan Guo Wei}{Oct.~2024 -- Jan.~2025}
\begin{itemize}
  \item Investigated whether LLMs can reason over non-text modalities without additional training, by treating modality tokens as in-context representations. The proposed In-Context Representation Learning (ICRL) framework, tested on the molecular domain, is to our knowledge the first training-free framework of its kind. Published at NeurIPS 2025 (first author).
\end{itemize}

\row{Shanghai Artificial Intelligence Laboratory}{Shanghai, China}
\subrow{Trainee Researcher.\, Focus: evaluation of large generative models.}{Oct.~2023 -- Oct.~2024}
\begin{itemize}
  \item Designed the T2VHE protocol, a standardized human-evaluation pipeline for text-to-video models that reduces evaluation cost by $\sim$50\% while preserving annotation quality. Published at NeurIPS 2024 (first author).
  \item Co-built ConvBench, a hierarchical multi-turn benchmark for LVLMs that disentangles perception, reasoning, and creation. Published at NeurIPS 2024 Spotlight (co-author).
\end{itemize}

\row{NUS High Performance Computing for AI (HPC-AI) Lab}{Singapore}
\subrow{Research Intern.\, Focus: graph dataset condensation.}{Apr.~2023 -- Oct.~2024}
\begin{itemize}
  \item Investigated lossless graph dataset condensation. Proposed expanding-window trajectory matching with a curriculum-based supervision design that substantially closes the gap to full-graph training. Published at ICML 2024 (co-first author).
\end{itemize}

% ---------- PUBLICATIONS ----------
\section*{Publications}
{\footnotesize \textbf{Bold} = author of this CV;\; \textbf{*} = equal contribution;\; \textbf{\dag} = corresponding author.}\par
\vspace{2pt}

\begin{itemize}[leftmargin=1.4em, itemsep=4pt]

  \item \textbf{(AAAI'26)}~To Align or Not to Align: Strategic Multimodal Representation Alignment for Optimal Performance.\par
        Wanlong Fang, \me{Tianle Zhang}, Alvin Chan\corr.\par
        \textit{Proceedings of the 40th AAAI Conference on Artificial Intelligence, 2026.}

  \item \textbf{(NeurIPS'25)}~Can LLMs Reason Over Non-Text Modalities in a Training-Free Manner? A Case Study with In-Context Representation Learning.\par
        \me{Tianle Zhang}\eq, Wanlong Fang\eq, Jonathan Woo\eq, Paridhi Latawa, Deepak A. Subramanian, Alvin Chan\corr.\par
        \textit{Proceedings of the 39th Conference on Neural Information Processing Systems, 2025.}

  \item \textbf{(NeurIPS'24)}~Rethinking Human Evaluation Protocol for Text-to-Video Models: Enhancing Reliability, Reproducibility, and Practicality.\par
        \me{Tianle Zhang}, Langtian Ma, Yuchen Yan, Yuchen Zhang, Kai Wang, Yue Yang, Ziyao Guo, Wenqi Shao, Yang You, Yu Qiao, Ping Luo, Kaipeng Zhang\corr.\par
        \textit{Proceedings of the 38th Conference on Neural Information Processing Systems, 2024.}

  \item \textbf{(NeurIPS'24 Spotlight)}~ConvBench: A Multi-Turn Conversation Evaluation Benchmark with Hierarchical Capability for Large Vision-Language Models.\par
        Shuo Liu, Kaining Ying, Hao Zhang, Yue Yang, Yuqi Lin, \me{Tianle Zhang}, Chuanhao Li, Yu Qiao, Ping Luo, Wenqi Shao\corr, Kaipeng Zhang\corr.\par
        \textit{Proceedings of the 38th Conference on Neural Information Processing Systems, 2024.}

  \item \textbf{(ICML'24)}~Navigating Complexity: Toward Lossless Graph Condensation via Expanding Window Matching.\par
        Yuchen Zhang\eq, \me{Tianle Zhang}\eq, Kai Wang, Ziyao Guo, Yuxuan Liang, Xavier Bresson, Wei Jin, Yang You.\par
        \textit{Proceedings of the 41st International Conference on Machine Learning, 2024.}

  \item \textbf{(ICONIP'23)}~ADGCN: A Weakly Supervised Framework for Anomaly Detection in Social Networks.\par
        Zhixiang Shen\eq, \me{Tianle Zhang}\eq\corr, Haolan He\eq.\par
        \textit{International Conference on Neural Information Processing, 2023 (CCF-C).}

  \item \textbf{(IEEE EHB'22)}~Phase Detection Methods for Alzheimer's Disease Based on Improved Convolutional Neural Network Fine-tuning.\par
        \me{Tianle Zhang}, Rui Yin, Zhan Qu, Zhuo Chen.\par
        \textit{10th IEEE International Conference on e-Health and Bioengineering, 2022 (IEEE Xplore, SCOPUS, WoS).}

  \item \textbf{(IEEE ITME'22)}~Fast Classification of Alzheimer's Disease Based on Improved Convolutional Neural Network.\par
        Zhan Qu, \me{Tianle Zhang}, Rui Yin, Zhuo Chen, Teoh Teik Toe.\par
        \textit{IEEE International Conference on Information Technology in Medicine and Education, 2022 (IEEE Xplore).}

\end{itemize}

% ---------- AWARDS ----------
\section*{Awards and Scholarships}

\begin{enumerate}[leftmargin=1.6em]
  \item \textbf{NTU Research Scholarship (R-RSS)} --- \textit{Graduate College, NTU; full tuition + monthly stipend.} \hfill Sep.~2024
  \item \textbf{Excellent First-class Scholarship} --- \textit{University of Electronic Science and Technology of China.} \hfill Oct.~2023
  \item \textbf{Xiaomi Scholarship} --- \textit{Xiaomi Foundation $\times$ UESTC.} \hfill Oct.~2023
  \item \textbf{Tencent First-class Scholarship} --- \textit{Tencent Foundation $\times$ UESTC.} \hfill Oct.~2022
  \item \textbf{NTU Project Scholarship} --- \textit{Nanyang Technological University (undergraduate research visit).} \hfill Oct.~2022
  \item \textbf{Bronze Award, Chengdu 80 Fintech Product R\&D Competition} --- \textit{Chengdu Fintech Industry Innovation Center.} \hfill Oct.~2021
  \item \textbf{Special Prize, 7th UESTC ``Internet Plus'' Innovation \& Entrepreneurship Competition} --- \textit{UESTC.} \hfill Aug.~2021
  \item \textbf{Second Prize (Sichuan), 11th National College E-commerce Competition} --- \textit{National Ministry of Education.} \hfill Aug.~2021
  \item \textbf{Blockchain Financial Experiment Class Scholarship} --- \textit{UESTC.} \hfill Oct.~2020
\end{enumerate}

% ---------- ACADEMIC SERVICE ----------
\section*{Academic Service}

\noindent\textbf{Reviewer Recognition}
\begin{itemize}
  \item \textbf{ICML 2026 Gold Reviewer} --- awarded to top reviewers based on area-chair ratings of submitted reviews; includes complimentary conference registration.
\end{itemize}

\vspace{2pt}
\noindent\textbf{Conference Reviewer}
\begin{itemize}
  \item International Conference on Learning Representations (ICLR): 2026
  \item Conference on Neural Information Processing Systems (NeurIPS): 2025, 2026
  \item International Conference on Machine Learning (ICML): 2025, 2026
  \item ICLR 2025 Workshop on AI for Nucleic Acids (AI4NA)
\end{itemize}

\vspace{2pt}
\noindent\textbf{Workshop Committee}
\begin{itemize}
  \item ECCV 2024 Workshop on Dataset Distillation Challenge --- Committee Member.
\end{itemize}

% ---------- SKILLS ----------
\section*{Skills and Languages}

\noindent\textbf{Programming:} Python (primary), Solidity, C; experienced with DApp / smart-contract development.\par
\textbf{ML / DL:} PyTorch, Hugging Face, scikit-learn; multi-modal LLMs, diffusion models, GNNs, dataset distillation.\par
\textbf{Tools:} Git, Linux, CUDA, LaTeX, Weights \& Biases.\par
\textbf{Languages:} Chinese (Native); English (Fluent --- academic writing, presentation, NeurIPS poster experience).

\end{document}
